% Mise en page du cours en PDF.
% Ce fragment est insere par .opencode/scripts/export-pdf.sh dans le modele
% LaTeX de pandoc. Il ne modifie jamais le contenu des supports : il regle
% seulement le rendu, pour que la hierarchie des titres corresponde au gabarit
% du cours, que les tableaux a six colonnes et les lignes de commande longues
% tiennent dans la largeur d'une page A4, et que l'identite du prestataire
% apparaisse sur la couverture et sur chaque page.

% --- Couleurs ----------------------------------------------------------------
% Les trois valeurs sont relevees sur le site du prestataire : le bleu nuit
% porte la couverture et le filet des titres, le rouge la marque et les
% pastilles de LAB, le gris ardoise les textes secondaires.
\definecolor{fond}{HTML}{171B2B}
\definecolor{marque}{HTML}{EF4A5C}
\definecolor{ardoise}{HTML}{5A6072}
\definecolor{livrable}{HTML}{1F7A5C}

% --- Polices -----------------------------------------------------------------
% TeX Gyre Pagella pour le texte : elle porte tous les accents francais, les
% espaces insecables et les guillemets, sans configuration supplementaire.
% Les polices sont designees par leur nom de fichier et non par leur nom de
% famille : XeLaTeX ne connait alors que les fichiers de son propre arbre, et
% la conversion ne depend pas des polices installees sur la machine hote.
%
% L'option Ligatures=NoCommon supprime les ligatures de mots, ff, fi et ffi.
% Elles sont esthetiquement correctes en francais, mais elles empechent de
% retrouver un mot par la recherche dans le PDF : un formateur qui cherche
% "suffixe" ou "coefficient" ne trouve rien. Dans un support forme, la
% recherche compte plus que la ligature.
\defaultfontfeatures{Ligatures=NoCommon}
\usepackage{fontspec}
\setmainfont{texgyrepagella-regular.otf}[
  BoldFont=texgyrepagella-bold.otf,
  ItalicFont=texgyrepagella-italic.otf,
  BoldItalicFont=texgyrepagella-bolditalic.otf
]
\setsansfont{texgyreheros-regular.otf}[
  BoldFont=texgyreheros-bold.otf,
  ItalicFont=texgyreheros-italic.otf,
  BoldItalicFont=texgyreheros-bolditalic.otf
]
\setmonofont{texgyrecursor-regular.otf}[
  BoldFont=texgyrecursor-bold.otf,
  ItalicFont=texgyrecursor-italic.otf,
  BoldItalicFont=texgyrecursor-bolditalic.otf,
  Scale=MatchLowercase
]

% --- Identite du prestataire -------------------------------------------------
% Le logo est un fichier a fond transparent, de sorte qu'il doit toujours etre
% pose sur le bandeau bleu nuit. Son chemin absolu est fourni par le script
% d'export, qui le connait : un chemin relatif dependrait du repertoire depuis
% lequel la conversion est lancee.
%
% graphicx est charge ici, et non laisse au modele de pandoc : ce modele ne
% l'inclut que si le support contient une image, or nos supports n'en
% contiennent aucune. Sans cet ajout, la couverture ne serait pas composee.
\usepackage{graphicx}
\providecommand{\ascentlogo}{logo-ascent.png}
\providecommand{\ascentintitule}{Forensique réseau}
\providecommand{\ascentjour}{}
\providecommand{\ascentdate}{}
\providecommand{\ascenttheme}{}
\providecommand{\ascentorganisme}{DGFIP}
\providecommand{\ascentprestataire}{Ascent Formation}
\providecommand{\ascentlieu}{Paris}

% Une pastille coloree. Elle sert au logo dans l'en-tete comme aux reperes
% LAB et LIVRABLE devant les titres : c'est le meme objet graphique, donc
% elle se definit une seule fois.
\newcommand{\ascentpastille}[2]{%
  \begingroup
  \setlength{\fboxsep}{1.8pt}%
  \colorbox{#1}{\textcolor{white}{\sffamily\bfseries\scriptsize #2}}%
  \endgroup
}

% Le bandeau de couverture : la largeur du texte, en bleu nuit, avec le logo
% blanc aligne a gauche et a distance reguliere du bord.
%
% La marge est posee a l'interieur de la boite de largeur fixe. Posee dehors,
% elle ajouterait sa longueur a celle de la boite et ferait deborder le
% bandeau de la page : c'est une erreur qui ne se voit pas sur le modele.
\newcommand{\ascentbandeau}{%
  \par\noindent
  \begingroup
  \setlength{\fboxsep}{0pt}%
  \colorbox{fond}{%
    \parbox[c][24mm][c]{\textwidth}{%
      \hspace*{10mm}%
      \includegraphics[width=42mm]{\ascentlogo}}}%
  \endgroup
  \par
}

% --- Langue ------------------------------------------------------------------
% babel est deja charge par le modele LaTeX de pandoc, qui recoit l'option
% lang=fr. On ne le recharge pas ici : deux chargements de babel avec des
% options differentes provoquent une erreur d'options incompatibles. On se
% contente de regler le comportement de la typographie francaise.
\frenchsetup{AutoSpacePunctuation=false}

% --- Tableaux ----------------------------------------------------------------
% Les tableaux du cours comptent jusqu'a six colonnes et certaines cellules
% contiennent une commande entiere. A la taille normale, ils debordent de la
% largeur du texte. On les ramene donc a la taille des notes de bas de page, ce
% qui reste lisible a l'impression et tient sur une largeur A4.
\usepackage{etoolbox}
\AtBeginEnvironment{longtable}{\footnotesize\setlength{\tabcolsep}{3pt}}
\AtBeginEnvironment{table}{\footnotesize\setlength{\tabcolsep}{3pt}}
\renewcommand{\arraystretch}{1.15}

% --- Blocs de code ----------------------------------------------------------
% Les commandes Shell les plus longues depassent cent vingt caracteres. Sans
% ces options, elles sortent de la page ou forcent le document a s'etaler
% horizontalement. On autorise donc la coupure des lignes longues.
\usepackage{fvextra}
\fvset{
  breaklines=true,
  breakanywhere=true,
  fontsize=\footnotesize,
  frame=single,
  framerule=0.3pt,
  framesep=4pt,
  breakindent=0pt,
  breakautoindent=false
}

% --- Liens ------------------------------------------------------------------
% Les references citent des adresses web longues, placees dans des cellules de
% tableau. Sans xurl, une adresse ne peut se couper qu'a certains caracteres et
% deborde de la page. xurl autorise la coupure n'importe ou.
\usepackage{xurl}
\Urlmuskip=0mu plus 1mu
\urlstyle{same}

% --- Titres ------------------------------------------------------------------
% Le script demande a pandoc de descendre les titres d'un cran, si bien que le
% titre du jour tient lieu de titre du document, un titre de niveau deux du
% gabarit devient une section, et une rubrique de chapitre devient une
% sous-section. La hierarchie du rendu est donc exactement celle du gabarit,
% sans avoir a retoucher les supports.
\usepackage{titlesec}
\usepackage{needspace}

% Un titre de section occupe une page entiere lorsqu'il commence par
% « Chapitre » : c'est ce qui donne au support sa lecture de livre. Les autres
% sections, liminaires ou finales, s'enchainent sans coupure. La rubrique qui
% nomme le livrable d'un laboratoire porte une pastille verte, elle aussi
% posee par le script d'export.
%
% Aucun de ces reperes ne peut etre pose ici, et l'echec des trois approches
% suivantes merite d'etre consigne.
%   - \renewcommand{\section}[2] : \section sert \everypar pour placer le
%     \label que Pandoc ecrit juste apres le titre. La redéfinir supprime ce
%     mecanisme, et l'on obtient « Argument of \label has an extra } ».
%   - \pretocmd{\section}{...{#2}...} : etoolbox ne transmet pas les
%     parametres de la commande corrigee. Le titre arrive sous la forme du
%     litteral « #2 », et le test echoue sans rien signaler.
%   - \IfBeginWith de xstring : cette commande consomme le jeton qui suit son
%     argument, et avale la meme accolade que \renewcommand.
% Reposer un compteur et des commandes numerotees, sans lire le titre, ne
% fonctionne pas non plus : le saut de page insere dans \section suffit a
% perdre le titre lui-meme et le paragraphe qui le suit.
%
% Le script d'export travaille donc le fichier .tex produit par Pandoc, juste
% avant la compilation : il insere le \clearpage devant les titres de chapitre
% et la pastille dans le titre du livrable. Le fragment se limite ensuite au
% style, ce qui est son role.
\renewcommand{\sectionmark}[1]{\markboth{\nouppercase{#1}}{}}

\titleformat{\section}
  {\needspace{5\baselineskip}\normalfont\sffamily\bfseries\LARGE\color{fond}}
  {}{0pt}{}[\vspace{2mm}{\color{marque}\titlerule[1.1pt]}]

\titleformat{\subsection}
  {\needspace{4\baselineskip}\normalfont\sffamily\bfseries\large\color{fond}}
  {}{0pt}{}

\titleformat{\subsubsection}
  {\needspace{3\baselineskip}\normalfont\sffamily\bfseries\normalsize\color{fond}}
  {}{0pt}{}

\titleformat{\paragraph}
  {\needspace{2\baselineskip}\normalfont\sffamily\bfseries\small\color{ardoise}}
  {}{0pt}{}

\titlespacing*{\section}{0pt}{0pt plus 1ex}{14pt}
\titlespacing*{\subsection}{0pt}{16pt plus 2ex}{5pt}
\titlespacing*{\subsubsection}{0pt}{11pt plus 1ex}{3pt}
\titlespacing*{\paragraph}{0pt}{9pt plus 1ex}{2pt}
\setlength{\parskip}{0.35em}

% --- Table des matieres ------------------------------------------------------
% Les quatre chapitres d'une journee et leurs rubriques doivent se voir d'un
% coup d'oeil, sans les sous-rubriques de LAB qui allongeraient la liste.
\setcounter{tocdepth}{2}
\renewcommand{\contentsname}{Sommaire}

% --- En-tete et pied de page -------------------------------------------------
% A gauche, la section en cours : un support de formation se cherche, et le
% lecteur doit savoir ou il est sans feuilleter. A droite, le logo du
% prestataire sur son fond bleu nuit, present sur chaque page.
\usepackage{fancyhdr}
\pagestyle{fancy}
\fancyhf{}
\fancyhead[L]{\footnotesize\color{ardoise}\leftmark}
\fancyhead[R]{%
  \begingroup
  \setlength{\fboxsep}{1.2pt}%
  \colorbox{fond}{\includegraphics[width=15mm]{\ascentlogo}}%
  \endgroup
}
\fancyfoot[C]{\small\color{fond}\thepage}
\renewcommand{\headrulewidth}{0.3pt}
\renewcommand{\footrulewidth}{0pt}
\setlength{\headheight}{18pt}

% --- Page de garde -----------------------------------------------------------
% L'intitule et la journee sont deux lignes distinctes : le script les extrait
% du premier titre du support, dont il ne garde que la partie qui precede le
% « JOUR X ». Afficher le titre de Pandoc en entier,echoiterait « JOUR 1 » deux
% fois de suite sur la couverture.
% Le bandeau, l'intitule, la journee, le theme, puis les mentions de
% l'organisme et du prestataire composent la couverture.
\makeatletter
\renewcommand{\maketitle}{%
  \thispagestyle{empty}%
  \ascentbandeau
  \vfill
  {\color{fond}\sffamily\bfseries\fontsize{28}{32}\selectfont
   \ascentintitule\par}%
  \vspace{5mm}%
  {\color{marque}\sffamily\bfseries\fontsize{20}{24}\selectfont
   \ascentjour\par}%
  \vspace{3mm}%
  {\color{ardoise}\sffamily\fontsize{12}{16}\selectfont \ascenttheme\par}%
  \vspace{3mm}%
  {\color{marque}\rule{45mm}{1.2pt}\par}%
  \vfill
  {\color{fond}\sffamily\fontsize{12}{16}\selectfont \ascentorganisme\par}%
  \vspace{2mm}%
  {\color{ardoise}\sffamily\fontsize{10}{14}\selectfont
   \ascentprestataire \textemdash{} \ascentlieu \textemdash{} \ascentdate\par}%
  \vfill
  \clearpage
}
\makeatother
